\uuid{DIEv}
\exo7id{2738}
\titre{exo7 2738}
\auteur{tumpach}
\organisation{exo7}
\datecreate{2009-10-25}
\isIndication{false}
\isCorrection{false}
\chapitre{Déterminant, système linéaire}
\sousChapitre{Système linéaire, rang}
\cpge{mpsi-mp2i.calcul-matriciel-systemes.c, pcsi.calcul-matriciel-systemes-lineaires.c, ptsi.calcul-matriciel-systemes-lineaires.c, mpsi-mp2i.matrices-applications-lineaires.c, pcsi.matrices-applications-lineaires.d, ptsi.matrices-applications-lineaires.d, bcpst1.systemes-lineaires.a, bcpst1.systemes-lineaires.c}
\module{Algèbre}
\niveau{L2}
\difficulte{}

\contenu{
\texte{
Soient $a$, $b$, et $c$ trois nombres r\'eels.
}
\begin{enumerate}
    \item \question{Quelle relation doivent satisfaire les param\`etres $a$, $b$ et $c$ pour que le syst\`eme suivant ait au moins une solution~?
$$
\mathcal{S}_{abc}~: \left\{\begin{array}{ccccccc} 
x & + & 2y & - & 3z & = & a\\
2x & + & 6y & - & 11z & = & b\\
x & - & 2y & + & 7z & = & c
\end{array}\right.
$$}
    \item \question{Est-ce que le syst\`eme $\mathcal{S}_{abc}$ peut avoir une unique solution~?}
\end{enumerate}
}
